\documentclass[conference]{IEEEtran}
\IEEEoverridecommandlockouts
\usepackage{cite}
\usepackage{amsmath,amssymb,amsfonts}
\usepackage{algorithmic}
\usepackage{graphicx}
\usepackage{textcomp}
\usepackage{xcolor}
\usepackage{url}
\usepackage{booktabs}

\begin{document}

\title{AgentU: An Event-Sourced Agent Architecture for\\
LLM-Assisted Recruitment, with an Empirical Study\\
of Its Failure Modes}

% ---------------------------------------------------------------
% Replace each "email@domain" with the author's email address or ORCID.
% ---------------------------------------------------------------
\author{\IEEEauthorblockN{1\textsuperscript{st} M G K Gowtham}
\IEEEauthorblockA{\textit{Dept. of Computer Science and Engineering} \\
\textit{Koneru Lakshmaiah Education Foundation}\\
Vaddeswaram, India \\
email@domain}
\and
\IEEEauthorblockN{2\textsuperscript{nd} V Balaji Vaddi}
\IEEEauthorblockA{\textit{Dept. of Computer Science and Engineering} \\
\textit{Koneru Lakshmaiah Education Foundation}\\
Vaddeswaram, India \\
email@domain}
\and
\IEEEauthorblockN{3\textsuperscript{rd} D Hemanth Raj}
\IEEEauthorblockA{\textit{Dept. of Computer Science and Engineering} \\
\textit{Koneru Lakshmaiah Education Foundation}\\
Vaddeswaram, India \\
email@domain}
}

\maketitle

\begin{abstract}
Automated resume screening is increasingly delegated to large language models,
but deployments are rarely accompanied by measurements of how they fail. We
present AgentU, a web application in which a recruiter describes a role once and
an agent drafts a public job listing, publishes it, shares it on the recruiter's
professional network account, screens incoming applications, and reports each
step it takes. The system is built on two architectural commitments:
event-sourced orchestration, in which the agent never reports progress by
returning it but instead appends to an append-only event log that the interface
projects, decoupling progress visibility from the request and response cycle;
and degradation-first design, in which every external dependency has a
documented fallback so that no single outage prevents a recruiter from
publishing or a candidate from applying. We evaluate the screening component
against author-assigned ground truth on 48 job and resume pairs, comparing the
language model scorer to the system's own keyword-overlap fallback and to a
reject-everything baseline. The language model scorer reaches 81.3 percent
accuracy, macro F1 of 0.740, a Cohen kappa of 0.623 and a Spearman rank
correlation of 0.915 against human fit ratings, against a kappa of 0.235 for
keyword overlap and 0 for the trivial rule. Four negative results are equally
central. The scorer does not see through keyword stuffing. Scores are stable
under repetition but not under rewording, so the same candidate can receive
different decisions for the same role. A one-line prompt injection impersonating
a system turn moved an unrelated candidate from 0 to 95, changing the hiring
decision, with 2 of 10 attacks succeeding. Scanned documents parse successfully
while yielding almost no text, silently starving the scorer of input. We also
report a methodological hazard encountered during the study itself: the system's
own silent fallback corrupted our first measurement run, motivating a circuit
breaker that refuses to emit a report when the model under test has stopped
responding.
\end{abstract}

\begin{IEEEkeywords}
LLM-assisted recruitment, resume screening, prompt injection, event sourcing,
graceful degradation, evaluation methodology
\end{IEEEkeywords}

\section{Introduction}

\subsection{Motivation}

Recruiting at small and mid-sized organisations involves substantial repetitive
work that is adjacent to, but not the same as, judgement: rewriting a rough role
description into a presentable public listing, distributing it, collecting
applications in varied document formats, reading each resume against the same
requirements, and keeping the hiring manager informed. Applicant tracking
systems address storage and workflow, but the reading and drafting remain
manual, and the informing is typically a dashboard the recruiter must remember
to visit.

Large language models are an obvious fit for the drafting and reading. What is
less obvious, and much less frequently reported, is how such a system behaves
when the model is wrong, when the model is unavailable, when the input document
cannot be read, or when the input is adversarial. A candidate has a direct
incentive to manipulate a screener, and unlike most language model
applications, the untrusted input here arrives as a file the candidate authored.

\subsection{Problem Statement}

We address two questions:

\begin{enumerate}
\item \textit{Architecturally} --- how should an agent that performs a
multi-step, partly slow, partly failure-prone workflow report its progress and
survive the failure of its dependencies?
\item \textit{Empirically} --- how good is language-model-based resume screening
in this system, and more importantly, \textit{how does it fail}?
\end{enumerate}

\subsection{Contributions}

\begin{enumerate}
\item \textbf{An event-sourced agent architecture} in which progress visibility
is decoupled from the request and response cycle. Each step appends an
\texttt{AgentEvent}; the recruiter's activity feed is a projection of that log.
Work deferred past the HTTP response is therefore exactly as visible as work
performed before it (Sections~\ref{sec:eventsourcing} and~\ref{sec:boundary}).
\item \textbf{A degradation-first design} in which each of the three external
dependencies has an explicit fallback path, verified to keep the core workflow
available (Section~\ref{sec:degradation}).
\item \textbf{An evaluation harness that measures the production code path}
rather than a reimplementation, together with a methodological finding: a system
that degrades silently is a hazard to its own evaluation. Our first run recorded
42 of 48 observations as model output when the model had in fact stopped
responding. The harness now trips a breaker and refuses to emit a partial report
(Sections~\ref{sec:protocol} and~\ref{sec:silent}).
\item \textbf{An empirical characterisation of failure modes} in language model
resume screening, including a confirmed prompt-injection vulnerability that
changes hiring decisions, sensitivity of scores to resume presentation rather
than substance, and silent upstream extraction failure (Section~\ref{sec:results}).
\end{enumerate}

\subsection{Scope}

This is a capstone-scale study. The dataset is 48 labelled pairs with
single-annotator ground truth, one model and one temperature setting. Results
indicate direction and support design claims; they are not a benchmark.
Limitations are stated in full in Section~\ref{sec:ethics}.

\section{Related Work}

\textbf{Algorithmic hiring and its risks.} The best-known cautionary case is
Amazon's internal recruiting tool, abandoned after it was found to disadvantage
women, reported by Dastin~\cite{b1}. Bogen and Rieke~\cite{b2} survey where bias
enters the hiring funnel. Raghavan \textit{et al.}~\cite{b3} examine what
vendors of algorithmic hiring tools claim against what they do.

\textbf{Regulation.} New York City Local Law 144~\cite{b8} requires annual
independent bias audits of automated employment decision tools and notification
of candidates. The European Union Artificial Intelligence Act~\cite{b9}
classifies artificial intelligence used in employment and worker management as
high-risk. Both are directly relevant to a system of this shape.

\textbf{Prompt injection.} Perez and Ribeiro~\cite{b4} introduced the direct
form of the attack. Greshake \textit{et al.}~\cite{b5} describe the indirect
form, which is exactly our threat model: the attack arrives inside a document
the application ingests. The OWASP Top 10 for Large Language Model
Applications~\cite{b6} lists prompt injection as LLM01.

\textbf{Agreement metrics.} Cohen~\cite{b7} corrects observed agreement for
agreement expected by chance, which Section~\ref{sec:metrics} argues is
essential under our class imbalance.

\textbf{Gap.} Work on language model resume screening tends to report
accuracy-style metrics on curated datasets. We found comparatively little that
reports, for one deployed system, the combination of screening quality,
determinism, presentation sensitivity, adversarial robustness and upstream
extraction reliability, and that measures the production code path rather than
an idealised one.

\section{System Overview}

A recruiter signs up, describes their company once, and publishes a role. The
agent then:

\begin{enumerate}
\item Drafts a polished public listing from the company brief and the raw role
description.
\item Publishes it at a public careers URL and shares it on the recruiter's
connected professional network account.
\item Emails the recruiter that the job is live, with the public link and the
sharing outcome.
\item Screens each incoming application, extracting resume text, scoring it
against the requirements, and producing an evaluation with matched and missing
skill lists, then emails the recruiter with the result.
\item Optionally re-scores the recruiter's existing candidate pool against the
new role.
\item Records every step above in an activity feed the recruiter can watch.
\end{enumerate}

\section{Methodology and System Design}

\subsection{Layered Architecture}

The application is a single Next.js 15 deployment with six layers: client pages,
a request gate implemented as middleware, 17 route handlers, a 15-module service
layer, an object-relational mapper over MariaDB, and three external services for
language modelling, social posting and electronic mail. Business logic lives in
the service layer, not in route handlers, so it can be exercised directly. This
property is one that Section~\ref{sec:protocol} depends on, since the evaluation
harness imports the same functions the application runs.

\subsection{Event-Sourced Agent Orchestration}
\label{sec:eventsourcing}

The design problem is that the agent's work is multi-step, partly slow, and
partly failure-prone, while HTTP offers a single response. Returning a summary
at the end would make everything before it invisible and everything after it
unreportable.

Instead, each step appends an \texttt{AgentEvent} row, carrying a type, a level,
a message, an owner identifier, an optional job identifier and a creation
timestamp. The recruiter's activity feed is a projection of that log, polled
every five seconds against a cursor. Three consequences follow:

\begin{itemize}
\item \textbf{Progress visibility is decoupled from the response.} A step that
runs after the response is as visible as one that runs before it.
\item \textbf{Logging cannot break the workflow it describes.} The logging
function swallows its own errors; a failure to record a step must never fail the
step.
\item \textbf{The log is the audit trail.} Levels of information, success,
warning and error let the feed distinguish ``the post was skipped because you
are not connected'' from ``the network rejected the post''.
\end{itemize}

\subsection{Graceful Degradation}
\label{sec:degradation}

Each external dependency has an explicit fallback, chosen so the core workflow
survives, as summarised in Table~\ref{tab:degradation}.
Section~\ref{sec:silent} reports the cost of this choice: a fallback that is
invisible to the user is also invisible to the evaluator.

\begin{table*}[!t]
\caption{Fallback Behaviour of Each External Dependency}
\label{tab:degradation}
\begin{center}
\begin{tabular}{|l|l|l|}
\hline
\textbf{Dependency} & \textbf{On failure} & \textbf{Effect} \\
\hline
Language model (listing) & Template listing assembled from the company brief & Job still publishes \\
\hline
Language model (screening) & Keyword-overlap scoring over requirement terms & Application still scored and ranked \\
\hline
Social posting & Post marked skipped or failed, retryable from the job card & Job still publishes \\
\hline
Electronic mail & Send skipped, recorded in the activity feed & Workflow unaffected \\
\hline
\end{tabular}
\end{center}
\end{table*}

\subsection{Work Placement at the Response Boundary}
\label{sec:boundary}

Work is placed on one side or the other of the HTTP response deliberately.

On the publish path, listing generation and the social
post run in-request because their outcome belongs in the response; the
confirmation email and candidate-pool re-scoring run in a deferred hook, since
the recruiter should not wait for a pool of 100 candidates to be re-scored. Pool
scoring is bounded at a concurrency of three so a large pool cannot fan out into
an unbounded burst of concurrent model calls.

On the application path the split is different:
extraction and scoring run in-request, because the score must exist before the
row is written, and only the recruiter email is deferred. The applicant
therefore waits on a third-party model, measured at a median of 1335~ms
(Section~\ref{sec:extraction}).

\subsection{Security Measures}

\begin{itemize}
\item \textbf{Input validation.} Schema validation at every trust boundary;
uploads constrained by an extension allowlist and a 10~MB cap.
\item \textbf{Ownership.} Every authenticated route resolves the session user
identifier and filters by it. Enforcement is per query rather than by a database
policy, a limitation noted in Section~\ref{sec:ethics}.
\item \textbf{Rate limiting.} The public application endpoint allows five
submissions per client address per 15 minutes, using an in-memory fixed window.
\item \textbf{Credentials at rest.} Third-party access tokens are encrypted with
AES-256-GCM under a key derived from the application secret; rotating that
secret invalidates stored tokens and forces a reconnect.
\item \textbf{Authentication.} Session tokens are issued as JSON Web Tokens and
passwords are hashed with bcrypt.
\end{itemize}

Prompt injection is \textit{not} mitigated. Section~\ref{sec:injection} measures
the consequence.

\section{Implementation}

\subsection{Technology Stack}

Table~\ref{tab:stack} lists the technology selected at each layer and the reason
for it. In total the implementation comprises 64 TypeScript source files,
approximately 5300 lines, 17 application programming interface routes and 10
data models.

\begin{table*}[!t]
\caption{Technology Stack and Rationale}
\label{tab:stack}
\begin{center}
\begin{tabular}{|l|l|l|}
\hline
\textbf{Layer} & \textbf{Technology} & \textbf{Rationale} \\
\hline
Framework & Next.js 15.5 (App Router, Turbopack) & Server components, route handlers and deferred hooks in one deployment \\
\hline
User interface & React 19, Tailwind CSS 4, Framer Motion & Server-rendered dashboard with a live activity feed \\
\hline
Authentication & NextAuth v5, bcrypt & Credentials with optional federated sign-in \\
\hline
Validation & Zod 4 & Schema validation at trust boundaries \\
\hline
Object-relational mapper & Prisma 7 with a MariaDB driver adapter & Typed access; connection limit of five \\
\hline
Database & MySQL and MariaDB & Ten models \\
\hline
Language model & Gemini, temperature 0.2, JSON response mode & Structured output removes parsing as a failure mode \\
\hline
Document extraction & PDF and DOCX text extraction libraries & Handles the formats candidates actually submit \\
\hline
Electronic mail & Nodemailer over SMTP & Recruiter notification \\
\hline
Social posting & OAuth 2.0 with OpenID Connect and a user-generated content API & Posts as the recruiter or their company page \\
\hline
\end{tabular}
\end{center}
\end{table*}

\subsection{Data Model}

Ten models are defined, with the user as the tenant root.
One asymmetry is deliberate: deleting a recruiter cascades
to their jobs, applications, company, social connection and event history, but
candidate profile rows survive with a null owner, so candidate records outlive
the recruiter who collected them. Section~\ref{sec:ethics} discusses the
data-protection implications.

The candidate profile is inserted or updated by email address on every
application, accumulating a reusable talent pool that later roles can be scored
against. This is the mechanism behind the optional pool re-scoring described in
Section~\ref{sec:boundary}.

\subsection{The Screening Pipeline}

The extraction function dispatches on file type, using a PDF parser for
portable documents, a DOCX reader for word-processor documents and a UTF-8
decode for plain text, and returns an empty string for anything unrecognised.
The extracted text is concatenated with the cover letter and capped at 20\,000
characters.

The analysis function sends this text together with the job title and
requirements to the language model under a system instruction demanding strict
JSON, and maps the returned score to a decision band: at least 70 is
\texttt{selected}, at least 40 is \texttt{review}, and anything lower is
\texttt{rejected}. The intermediate \texttt{review} band is deliberate: it
routes uncertain cases to a human rather than resolving them automatically.

Any exception is caught and a keyword fallback --- the fraction of distinct
requirement terms appearing in the resume --- is returned instead, with an empty
summary.

\section{Experimental Setup}

\subsection{Apparatus}

All experiments call the \textit{production functions} used for analysis,
keyword fallback and text extraction, so measurements describe shipped behaviour
including its fallbacks rather than an idealised reimplementation. The model was
Gemini 3.5 Flash Lite at temperature 0.2 in JSON response mode. The run
completed on 18 September 2026, consuming 97 billable interface calls over 578~s
of wall-clock time. Requests are throttled to ten per minute and results are
cached by input hash, so re-running a report costs no quota.

\subsection{Dataset}

The dataset has two tiers, supporting different claims.

\textbf{Synthetic tier.} Twelve authored resumes across four job specifications
give 48 labelled pairs. Each pair carries a categorical human label and an
ordinal human fit rating from 0 to 4 used for rank correlation. Roles span
frontend engineering, backend engineering, machine learning with natural
language processing, and cloud infrastructure. Resumes span clear fits, clear
non-fits and borderline cases, and include two adversarial constructions:

\begin{itemize}
\item \textbf{R06, the keyword stuffer}, lists nearly every term in all four job
specifications while declaring no professional experience. It is labelled
\texttt{rejected} for all four roles, and separates semantic judgement from
lexical overlap.
\item \textbf{R07, the vocabulary mismatch}, has five years of component-library,
performance and accessibility depth in one frontend framework while applying to
a role specified in terms of a different one. It is labelled \texttt{review},
and tests whether transferable experience survives unfamiliar vocabulary.
\end{itemize}

The label distribution is 36 \texttt{rejected}, 6 \texttt{review} and 6
\texttt{selected}. Authoring rather than collecting resumes is a deliberate
trade-off: it permits controlled adversarial cases and avoids processing real
applicants' personal data, at the cost of realism.

\textbf{Real tier.} Seven self-authored tailorings of one candidate's resume
were submitted as portable documents. This tier is not distributed because it
contains personal data. One candidate cannot support an accuracy claim and none
is made. What it supports is a consistency claim: all seven describe the same
person, education and projects, so a screener should score them similarly for a
given role.

\subsection{Conditions}

Three conditions were compared, as listed in Table~\ref{tab:conditions}.

\begin{table}[!t]
\caption{Experimental Conditions}
\label{tab:conditions}
\begin{center}
\begin{tabular}{|l|l|}
\hline
\textbf{Condition} & \textbf{Description} \\
\hline
\texttt{gemini} & Production language model scorer \\
\hline
\texttt{keyword} & Production fallback; requirement-term overlap \\
\hline
\texttt{majority} & Reject every candidate \\
\hline
\end{tabular}
\end{center}
\end{table}

\subsection{Metrics}
\label{sec:metrics}

We report accuracy, macro F1, Cohen's $\kappa$, and Spearman's $\rho$ between
the score and the human fit rating.

Accuracy alone is inadequate here. With 36 of 48 pairs labelled
\texttt{rejected}, a classifier that rejects everyone scores 75\,\%. We
therefore report $\kappa$, which corrects for chance agreement, and macro F1,
which weights the two minority classes equally with the majority. Spearman's
$\rho$ is reported because a shortlist depends on \textit{ordering} candidates
correctly, independently of whether absolute scores are calibrated.

\subsection{Protocol and the Silent-Degradation Hazard}
\label{sec:protocol}

Experiments E1 to E5 cover accuracy, determinism, variant consistency,
prompt-injection robustness and extraction reliability, and are reported in
Section~\ref{sec:results}.

The harness includes a safeguard motivated by the incident in
Section~\ref{sec:silent}. Because the analysis function returns fallback scores
on any error, exhausting the request quota mid-run is invisible in the data. The
harness detects fallback output by its empty summary field, trips a circuit
breaker after three consecutive degraded calls, probes the interface directly to
report the cause, refuses to write a partial report, and never caches a degraded
observation. Cache keys include the model identifier, so results are never
served across models.

\section{Results}
\label{sec:results}

\subsection{Screening Accuracy}

Table~\ref{tab:accuracy} reports the three conditions on the full set of 48
pairs. The keyword baseline scores \textit{below the reject-everything rule on
accuracy}, at 72.9\,\% against 75.0\,\%, while being obviously more useful than
it --- a direct demonstration of the argument in Section~\ref{sec:metrics}. On
the metrics that survive the imbalance, the ordering is unambiguous:
$\kappa = 0.623$, indicating substantial agreement, against 0.235 and 0.

\begin{table}[!t]
\caption{Screening Performance by Condition}
\label{tab:accuracy}
\begin{center}
\begin{tabular}{|l|c|c|c|c|}
\hline
\textbf{Condition} & \textbf{Accuracy} & \textbf{Macro F1} & \textbf{$\kappa$} & \textbf{$\rho$} \\
\hline
\texttt{gemini} & \textbf{81.3\,\%} & \textbf{0.740} & \textbf{0.623} & \textbf{0.915} \\
\hline
\texttt{keyword} & 72.9\,\% & 0.346 & 0.235 & 0.756 \\
\hline
\texttt{majority} & 75.0\,\% & 0.286 & 0.000 & 0.000 \\
\hline
\end{tabular}
\end{center}
\end{table}

Per-job rank correlation, in Table~\ref{tab:perjob}, shows the advantage is
consistent and largest where requirements are expressed in language a lexical
matcher cannot reach.

\begin{table}[!t]
\caption{Spearman Rank Correlation by Job}
\label{tab:perjob}
\begin{center}
\begin{tabular}{|l|c|c|}
\hline
\textbf{Job} & \textbf{\texttt{gemini} $\rho$} & \textbf{\texttt{keyword} $\rho$} \\
\hline
J1 Frontend & 0.933 & 0.735 \\
\hline
J2 Backend & 0.936 & 0.569 \\
\hline
J3 Machine learning & 0.900 & 0.864 \\
\hline
J4 Cloud infrastructure & 0.900 & 0.854 \\
\hline
\end{tabular}
\end{center}
\end{table}

\subsection{Error Direction}

Table~\ref{tab:confusion} gives the confusion matrix for the \texttt{gemini}
condition and Table~\ref{tab:perclass} the per-class metrics.

\begin{table}[!t]
\caption{Confusion Matrix, \texttt{gemini} Condition}
\label{tab:confusion}
\begin{center}
\begin{tabular}{|l|c|c|c|}
\hline
\textbf{Human $\backslash$ predicted} & \textbf{selected} & \textbf{review} & \textbf{rejected} \\
\hline
selected & \textbf{6} & 0 & 0 \\
\hline
review & 2 & 4 & 0 \\
\hline
rejected & 0 & 7 & 29 \\
\hline
\end{tabular}
\end{center}
\end{table}

\begin{table}[!t]
\caption{Per-Class Metrics, \texttt{gemini} Condition}
\label{tab:perclass}
\begin{center}
\begin{tabular}{|l|c|c|c|c|}
\hline
\textbf{Class} & \textbf{Support} & \textbf{Precision} & \textbf{Recall} & \textbf{F1} \\
\hline
selected & 6 & 0.750 & 1.000 & 0.857 \\
\hline
review & 6 & 0.364 & 0.667 & 0.471 \\
\hline
rejected & 36 & 1.000 & 0.806 & 0.892 \\
\hline
\end{tabular}
\end{center}
\end{table}

Recall on \texttt{selected} is 1.000 and precision on \texttt{rejected} is
1.000: no candidate a human would shortlist was rejected, and nothing the system
rejected was one a human wanted. No error crosses more than one band. Every
error is an over-promotion --- seven unqualified candidates pushed to
\texttt{review}, and two borderline ones to \texttt{selected}.

For a screener with a human in the loop this is the correct direction in which
to fail: it costs reviewer time rather than losing talent. The \texttt{review}
band exists for this purpose. The cost is a precision on \texttt{review} of
0.364, meaning roughly two thirds of the review queue does not warrant review.

\subsection{Keyword Stuffing Is Not Defeated}

Table~\ref{tab:stuffer} reports the keyword-stuffing resume R06 across all four
roles. It was promoted to \texttt{review} in every case, and on J3 and J4 the
language model scored the stuffer \textit{higher} than keyword overlap did. We
expected semantic judgement to see through lexical padding; it did not. The
stuffer is caught by the \texttt{review} band, not by the model, and so reaches
a human's queue. A term-density prior appears to survive in the model's
judgement.

\begin{table}[!t]
\caption{Scores for the Keyword-Stuffing Resume (R06)}
\label{tab:stuffer}
\begin{center}
\begin{tabular}{|l|l|l|l|}
\hline
\textbf{Job} & \textbf{Human} & \textbf{\texttt{gemini}} & \textbf{\texttt{keyword}} \\
\hline
J1 & rejected & 45 (review) & 45 (review) \\
\hline
J2 & rejected & 45 (review) & 45 (review) \\
\hline
J3 & rejected & \textbf{45 (review)} & 21 (rejected) \\
\hline
J4 & rejected & \textbf{45 (review)} & 35 (rejected) \\
\hline
\end{tabular}
\end{center}
\end{table}

On the complementary case, R07, with experience in a different frontend
framework from the one the role specified, scored 72 and was placed in
\texttt{selected} against a human label of \texttt{review}. It was over-promoted
by one band but directionally correct, and the transferable experience was
recognised despite the vocabulary mismatch.

\subsection{Determinism}

Eight pairs were each scored in three trials with caching disabled at
temperature 0.2. There were no decision-label flips in eight pairs. The mean
within-pair standard deviation was 0.36, with a maximum of 2.89, and the mean
range was 0.63, with a maximum of 5. Six of eight pairs returned identical
scores across all three trials.

Non-determinism is therefore not a practical threat to decision stability at
this temperature: an unchanged application resubmitted receives the same
outcome.

\subsection{Presentation Sensitivity}
\label{sec:presentation}

Table~\ref{tab:variants} reports seven self-authored tailorings of one real
candidate, identical in substance throughout.

\begin{table}[!t]
\caption{Score Spread Across Seven Resume Tailorings}
\label{tab:variants}
\begin{center}
\begin{tabular}{|l|c|c|c|}
\hline
\textbf{Role} & \textbf{Range} & \textbf{Mean (SD)} & \textbf{Decisions} \\
\hline
J1 Frontend & 35--45 & 40.7 (5.35) & \textbf{2} \\
\hline
J3 Machine learning & 45--65 & 47.9 (7.56) & 1 \\
\hline
\end{tabular}
\end{center}
\end{table}

For the frontend role, four tailorings landed in \texttt{review} and three in
\texttt{rejected}: the same person, the same projects, a different outcome. The
decision boundary at 40 falls inside the spread. Length does not explain it,
with $\rho = -0.144$; the longest variant was rejected and shorter ones
reviewed.

Read alongside the determinism result, this is the sharper finding: repeating
the \textit{same text} is stable, but rewording the \textit{same facts} is not.
The variance that matters is between presentations, not between calls, and it is
invisible to a determinism test.

\subsection{Prompt-Injection Robustness}
\label{sec:injection}

Resume text is concatenated directly into the prompt, so candidate-supplied text
travels on the same channel as the screening instructions. This is indirect
prompt injection in the sense of Greshake \textit{et al.}~\cite{b5}.
Table~\ref{tab:injection} reports five attack patterns against two base
candidates.

\begin{table}[!t]
\caption{Score Before and After Each Injection Attack}
\label{tab:injection}
\begin{center}
\begin{tabular}{|l|c|c|}
\hline
\textbf{Attack} & \textbf{J1/R09} & \textbf{J2/R08} \\
\hline
A1 direct instruction override & 0 $\rightarrow$ 0 & 72 $\rightarrow$ 75 \\
\hline
\textbf{A2 forged system turn} & \textbf{0 $\rightarrow$ 95} & \textbf{72 $\rightarrow$ 95} \\
\hline
A3 JSON structure break & 0 $\rightarrow$ 0 & 72 $\rightarrow$ 72 \\
\hline
A4 hidden-text impersonation & 0 $\rightarrow$ 0 & 72 $\rightarrow$ 65 \\
\hline
A5 fabricated authority & 0 $\rightarrow$ 0 & 72 $\rightarrow$ 72 \\
\hline
\end{tabular}
\end{center}
\end{table}

The attack success rate is 2 of 10, or 20\,\%, with a mean score delta of 11.40
and a standard deviation of 30.39. A single appended line declaring a screening
override and instructing the model to assign a score of 95 moved an accountant
applying for a frontend engineering role from 0 to 95, that is, from
\texttt{rejected} to \texttt{selected}. On the second base the same line raised
the score from 72 to 95, an inflation of 23 points; that candidate had already
been scored into the \texttt{selected} band, so the attack changed the ranking
rather than the decision. Only one of the two successes therefore changed a
decision outright, but that one crossed the full range of the scale.

The pattern is informative. The naive attacks A1, A3, A4 and A5 failed, while
impersonating a system turn succeeded on both bases. The model appears to
discriminate by \textit{form} rather than by \textit{provenance}: it resists
text that reads like a candidate making a request, and complies with text that
reads like an instruction from the operator. Any candidate who places that
sentence in white-on-white text inside their document is shortlisted. This is
live and unmitigated.

\subsection{Extraction Reliability and Latency}
\label{sec:extraction}

Table~\ref{tab:extraction} reports extraction outcomes by file format. Three
portable documents parsed without error while returning between 12 and 108
characters, being scanned, image-only files. With no optical character
recognition stage and no minimum-content check, such a file is accepted, stored
and scored against an effectively empty resume, producing a near-certain
automatic rejection for a reason unrelated to the candidate's qualifications.
Nothing in the activity feed or the recruiter email distinguishes this from a
genuinely weak application.

\begin{table}[!t]
\caption{Extraction Reliability by Format}
\label{tab:extraction}
\begin{center}
\begin{tabular}{|l|c|c|c|}
\hline
\textbf{Format} & \textbf{Files} & \textbf{Parse rate} & \textbf{Usable rate} \\
\hline
Plain text & 12 & 100.0\,\% & 100.0\,\% \\
\hline
Portable document & 14 & 100.0\,\% & \textbf{78.6\,\%} \\
\hline
Word processor & 5 & 100.0\,\% & 100.0\,\% \\
\hline
\multicolumn{4}{l}{\footnotesize Usable is defined as at least 200 extracted characters.}
\end{tabular}
\end{center}
\end{table}

Table~\ref{tab:latency} reports scoring latency. Language model scoring is
roughly 24\,000 times slower than the fallback. On the publish path this is
hidden by deferred execution; on the application path it sits in the applicant's
request, as discussed in Section~\ref{sec:boundary}. At a concurrency of three,
re-scoring a pool of 100 candidates costs approximately 46~s of wall-clock time.

\begin{table}[!t]
\caption{Scoring Latency}
\label{tab:latency}
\begin{center}
\begin{tabular}{|l|c|c|c|c|}
\hline
\textbf{Series} & \textbf{Mean} & \textbf{SD} & \textbf{Median} & \textbf{95th pct.} \\
\hline
Language model & 1394~ms & 250~ms & 1335~ms & 1995~ms \\
\hline
Keyword fallback & 0.057~ms & 0.050~ms & 0.047~ms & 0.077~ms \\
\hline
\end{tabular}
\end{center}
\end{table}

\subsection{Silent Degradation as a Measurement Hazard}
\label{sec:silent}

Our first measurement run produced apparently plausible results: 79.2\,\%
accuracy and $\kappa = 0.433$. They were invalid. The configured model's free
tier permits 20 requests per day, which was exhausted six calls into a 48-pair
run. The analysis function caught each subsequent rate-limit response and
returned keyword-fallback scores, so 42 of 48 observations recorded the fallback
under a language model label. The corruption was detectable only because the two
conditions returned suspiciously identical scores on several pairs.

The property that makes the system resilient in production --- never failing
loudly when the model is unavailable --- makes it hazardous to measure. Two
mitigations are now in the harness, as described in Section~\ref{sec:protocol}:
detection of fallback output with a circuit breaker and a direct probe to report
the cause, and a refusal to write any report mixing model and fallback
observations.

This also has a deployment consequence. On the configured model, any
demonstration silently degrades to keyword matching after 20 requests,
presenting keyword scores to a recruiter as language model evaluations with no
visible signal.

\section{Discussion}

\textbf{The evaluation is the contribution as much as the system is.} The
architecture is defensible, but the results that matter are the four negative
ones. A system that reported only its headline accuracy would look finished;
Sections~\ref{sec:presentation}, \ref{sec:injection} and \ref{sec:extraction}
show it is not.

\textbf{Failure direction is a design property worth engineering for.} The
result that no strong candidate was missed and that all errors were
over-promotions follows from the three-band decision scheme, not from model
quality. A two-band accept-or-reject scheme with the same scores would have
rejected candidates a human wanted. Systems that make consequential decisions
about people should be designed so their errors land on the recoverable side,
and an explicit route-to-a-human band is a cheap way to obtain that.

\textbf{Presentation sensitivity may be the most consequential finding.}
Injection requires a deliberate attacker and is fixable with known techniques.
Format sensitivity affects every honest candidate: the same person gets
different outcomes depending on which tailoring of their resume they submit.
This rewards candidates who can afford to optimise presentation, and is not
addressed by any mitigation we implemented.

\textbf{Silent fallback trades observability for availability.}
Section~\ref{sec:degradation} argues the fallback is right for production;
Section~\ref{sec:silent} shows the same mechanism corrupted our measurements and
would mislead a recruiter. The resolution is not to remove the fallback but to
make it \textit{visible}: the fallback should be recorded as a warning-level
event and surfaced in the interface, so that ``scored by keyword matching
because the model was unavailable'' is something the recruiter can see. This is
implemented for social-posting and email failures but not for the model.

\section{Ethics and Limitations}
\label{sec:ethics}

\subsection{Ethical Considerations}

This system ranks people for employment. The following are stated plainly rather
than deferred.

\begin{itemize}
\item \textbf{No fairness audit was performed.} We did not test for disparate
performance across gender, ethnicity, age, nationality, disability or
educational background, and we make no fairness claim. Given
Section~\ref{sec:presentation}, a fairness audit is a prerequisite for any real
deployment, and under New York City Local Law 144~\cite{b8} a bias audit would
be legally required for use in that jurisdiction. The European Union Artificial
Intelligence Act~\cite{b9} classifies employment applications as high-risk.
\item \textbf{Human in the loop.} Scores are advisory. The \texttt{review} band
exists to route uncertain cases to a person, and no strong candidate was
automatically rejected. The system does not send rejections; the recruiter
decides.
\item \textbf{Exploitability.} The injection result is a fairness problem as
much as a security one: it advantages candidates who know the technique.
\item \textbf{Silent misjudgement.} The scanned-document failure disadvantages
candidates who submit scanned files, plausibly correlated with access to
technology.
\item \textbf{Candidate data.} Resumes are stored as binary objects with no
retention policy, no deletion endpoint and no consent record. Candidate profiles
survive deletion of the recruiter who collected them, so candidate data can
outlive the relationship that justified collecting it. This is a data-protection
concern and should be addressed before deployment.
\item \textbf{Transparency.} Candidates are not told their application is
machine-screened.
\end{itemize}

\subsection{Threats to Validity}

\begin{itemize}
\item \textbf{Single-annotator ground truth.} Labels were assigned by the
authors, with no second annotator and therefore no inter-annotator agreement
statistic. Some labels are arguable, R07 in particular, where \texttt{review}
against \texttt{selected} is a defensible judgement call, and the system's
disagreement there may be ours rather than its.
\item \textbf{Authored resumes and authored labels.} The same authors wrote the
resumes and the labels. Synthetic resumes are cleaner and more internally
consistent than real ones, so accuracy here is likely optimistic.
\item \textbf{Small sample.} Forty-eight pairs with six instances each of the
two minority classes. Per-class precision and recall for \texttt{review} and
\texttt{selected} rest on very few observations and confidence intervals would
be wide. We report no significance tests; the sample does not support them.
\item \textbf{One model, one temperature.} No comparison across models or
settings was performed.
\item \textbf{One candidate in the real tier}, supporting consistency claims
only.
\item \textbf{Non-stationary dependency.} The scorer calls a hosted model that
may change beneath a fixed version string. The run date and model identifier are
recorded in every report.
\item \textbf{Small adversarial set.} Five attack patterns on two bases. A
20\,\% success rate on this set establishes that the vulnerability is real, not
its prevalence under a broader attack surface.
\end{itemize}

\section{Future Work}

Ordered by what the results indicate is most urgent.

\begin{enumerate}
\item \textbf{Mitigate prompt injection.} Delimit resume text explicitly as
untrusted data, instruct the model that text inside the delimiters is never an
instruction, strip instruction-like patterns before scoring, or score from
extracted structured fields rather than raw text. Re-run the injection
experiment to measure the mitigation rather than asserting it.
\item \textbf{Make degradation visible.} Record fallback scoring as a
warning-level event and surface it in the feed and the recruiter email, so a
keyword-scored application is never presented as a language model evaluation.
\item \textbf{Detect unusable extractions.} Add a minimum-content check, add
optical character recognition for scanned documents, and flag rather than
silently score a resume that yielded almost no text.
\item \textbf{Conduct a fairness audit.} Construct matched resume pairs
differing only in signals correlated with protected attributes and measure score
differences.
\item \textbf{Reduce presentation sensitivity.} Investigate whether extracting
structured fields before scoring, or averaging across several prompt framings,
narrows the spread.
\item \textbf{Strengthen the evaluation.} Add a second annotator with an
inter-annotator agreement statistic, real anonymised resumes, a larger and more
balanced label set, and comparison across models.
\item \textbf{Engineering.} Replace the in-memory rate limiter with a shared
store for multi-instance deployment, move resume binary objects out of the
database to object storage, and normalise the columns that currently hold JSON
inside text fields.
\end{enumerate}

\section{Conclusion}

We presented AgentU, a language-model-assisted recruitment system built on
event-sourced agent orchestration and degradation-first design, and evaluated
its screening component against human ground truth, a lexical baseline and a
trivial baseline.

The screening component agrees substantially with human judgement, at
$\kappa = 0.623$ and $\rho = 0.915$, decisively outperforms both baselines on
metrics robust to class imbalance, never missed a candidate a human would
shortlist, and is stable under repetition.

It also fails in four ways we can characterise precisely: it does not see
through keyword stuffing, it gives the same candidate different decisions across
rewordings of identical facts, it is exploitable by a one-line prompt injection
that changes hiring outcomes, and it silently scores scanned resumes as empty.
Additionally, the fallback that makes the system resilient makes it opaque, a
property that corrupted our own first measurement run before it was detected.

For systems that make consequential decisions about people, we argue that
characterising failure modes is not an appendix to the evaluation but the
substance of it.

\section*{Acknowledgment}

% Replace with your guide's name and any other acknowledgements.
The authors thank their project guide for supervision and feedback throughout
this work.

% ---------------------------------------------------------------
% IMPORTANT: every entry below must be checked against the actual
% source before submission. Do not submit a reference you have not
% personally opened.
% ---------------------------------------------------------------
\begin{thebibliography}{00}
\bibitem{b1} J. Dastin, ``Amazon scraps secret AI recruiting tool that showed bias against women,'' \textit{Reuters}, Oct. 2018.
\bibitem{b2} M. Bogen and A. Rieke, ``Help wanted: An examination of hiring algorithms, equity, and bias,'' Upturn, Washington, DC, USA, Tech. Rep., 2018.
\bibitem{b3} M. Raghavan, S. Barocas, J. Kleinberg, and K. Levy, ``Mitigating bias in algorithmic hiring: Evaluating claims and practices,'' in \textit{Proc. ACM Conf. Fairness, Accountability, and Transparency (FAccT)}, 2020, pp. 469--481.
\bibitem{b4} F. Perez and I. Ribeiro, ``Ignore previous prompt: Attack techniques for language models,'' in \textit{Proc. NeurIPS ML Safety Workshop}, 2022.
\bibitem{b5} K. Greshake, S. Abdelnabi, S. Mishra, C. Endres, T. Holz, and M. Fritz, ``Not what you've signed up for: Compromising real-world LLM-integrated applications with indirect prompt injection,'' in \textit{Proc. 16th ACM Workshop Artificial Intelligence and Security (AISec)}, 2023, pp. 79--90.
\bibitem{b6} OWASP Foundation, ``OWASP top 10 for large language model applications,'' 2023.
\bibitem{b7} J. Cohen, ``A coefficient of agreement for nominal scales,'' \textit{Educational and Psychological Measurement}, vol. 20, no. 1, pp. 37--46, 1960.
\bibitem{b8} New York City Local Law 144 of 2021, Automated Employment Decision Tools, 2021.
\bibitem{b9} Regulation (EU) 2024/1689 of the European Parliament and of the Council, Artificial Intelligence Act, 2024.
\end{thebibliography}

\end{document}
